% =====================================================================
\section{TinyML: нейросеть внутри микроконтроллера}\label{les:tinyml}
% =====================================================================
\lessonintro
  {пройти весь путь TinyML --- от данных до работающей нейросети на плате --- и измерить качество, время и память}
  {урок 4 (ШИМ), урок 6 (АЦП), урок 26 (акселерометр), урок 37 (обучение, инференс, честная проверка)}
  {сеть, которая повторяет синус и управляет яркостью светодиода, и классификатор жестов по акселерометру}
  {плата, светодиод, MPU6050 (для опыта 2); компьютер с Python и NumPy; PlatformIO по желанию (урок 41)}

\subsection{Сначала: что такое TinyML}

В уроке 37 мы научили программу различать два положения ручки, запомнив «центры» двух классов.
Это уже машинное обучение: правило получено не от программиста, а из примеров. \textbf{TinyML} ---
направление, в котором такие модели, в том числе нейронные сети, запускают на микроконтроллерах,
где нет операционной системы, гигабайт памяти и розетки. «Tiny» (крошечный) относится ко всему
сразу: к модели, к памяти, к потреблению энергии.

Зачем это нужно, если есть облако? Сравните:

\begin{center}\small
\begin{tabularx}{\textwidth}{@{}l X X@{}}
\toprule
 & Модель в облаке & Модель на микроконтроллере \\
\midrule
Задержка & сотни миллисекунд и зависит от сети & единицы или десятки миллисекунд, стабильно \\
Без интернета & не работает & работает \\
Приватность & звук и кадры уходят с устройства & наружу уходит только решение («жест: тряска») \\
Энергия & радиомодуль включён постоянно & радио спит, пока нет события (урок 9) \\
Размер модели & сколько угодно & от сотен байт до сотен килобайт \\
Сложность задачи & речь, текст, сложные изображения & жесты, шаги, вибрация, ключевое слово, простая картинка \\
\bottomrule
\end{tabularx}
\end{center}

Типичные задачи TinyML: «проснись по слову», распознавание жестов и активности по акселерометру,
обнаружение износа подшипника по вибрации, подсчёт шагов, определение, что в кадре есть человек.
Во всех случаях данные приходят от датчика \emph{на самом устройстве}, а решение нужно сразу.

\subsection{Когда нейросеть не нужна}

Нейросеть --- не волшебство, а ещё один способ написать функцию «признаки $\to$ решение».
Если задачу решает порог (температура выше 30 градусов) или простая формула, нейросеть лишь
усложнит проект: нужны данные, обучение, проверка и больше памяти. Нейросеть оправданна, когда
правило трудно записать вручную: чем жест «тряска» отличается от жеста «встряхнул и положил»
словами не скажешь, а примеры собрать легко. Всегда сравнивайте с простым решением,
как мы делали в уроке 37.

\subsection{Рабочий цикл TinyML}

\begin{figure}[H]
\centering
\begin{tikzpicture}[font=\sffamily\scriptsize,
  f/.style={blk,minimum width=23mm,minimum height=12mm,align=center}]
  \node[f,fill=blkper] (d) at (0,0) {\textbf{1. Данные}\\ записи с датчика\\ с метками классов};
  \node[f,fill=blkper] (p) at (3.9,0) {\textbf{2. Признаки}\\ окно, среднее,\\ разброс};
  \node[f,fill=blkcpu] (t) at (7.8,0) {\textbf{3. Обучение}\\ на компьютере,\\ много раз};
  \node[f,fill=blkmem] (q) at (11.7,0) {\textbf{4. Квантование}\\ веса float $\to$ int8};
  \node[f,fill=blkmem] (c) at (11.7,-2.3) {\textbf{5. В код}\\ массив чисел\\ в прошивке};
  \node[f,fill=blkrf] (i) at (7.8,-2.3) {\textbf{6. Инференс}\\ на плате,\\ каждое окно};
  \node[f,fill=blkrtc] (m) at (3.9,-2.3) {\textbf{7. Измерение}\\ точность, время,\\ память};
  \draw[arr] (d) -- (p); \draw[arr] (p) -- (t); \draw[arr] (t) -- (q);
  \draw[arr] (q) -- (c); \draw[arr] (c) -- (i); \draw[arr] (i) -- (m);
  \draw[arr,dashed] (m.south) -- ++(0,-0.6) -| node[pos=0.25,below]{плохо: нужно больше данных} (d.south);
\end{tikzpicture}
\caption{Цикл TinyML. Обучение идёт на компьютере, на плате остаётся только прямой расчёт.
Результат измерений возвращает нас к данным}
\end{figure}

Главное отличие от обычной программы: ошибка чаще прячется в \emph{данных}, а не в коде.
Поэтому шаги 1 и 7 важнее всего остального.

\subsection{Нейрон и слой}

Искусственный нейрон делает три вещи: умножает каждый вход на свой вес, складывает и добавляет
смещение, затем пропускает результат через «функцию активации». В этой главе --- функция ReLU,
которая заменяет отрицательное число нулём:
\[
y = \operatorname{ReLU}(w_1x_1 + w_2x_2 + b),\qquad \operatorname{ReLU}(z)=\max(0,z).
\]
Пример. Пусть $x=(0{,}5;\,-1)$, веса $w=(2;\,1)$, смещение $b=0{,}3$. Тогда
$2\cdot0{,}5 + 1\cdot(-1) + 0{,}3 = 0{,}3$, и $y = 0{,}3$. При $b=-0{,}5$ сумма равна $-0{,}5$,
ReLU даёт $y=0$: нейрон «не сработал».

\textbf{Слой} --- несколько таких нейронов, получающих одни и те же входы. Для слоя с $n_\text{in}$
входами и $n_\text{out}$ нейронами нужна таблица весов $n_\text{in}\times n_\text{out}$ и $n_\text{out}$
смещений. \textbf{Сеть} --- цепочка слоёв. Число параметров считают просто:
\[
N = \sum_{\text{слои}} (n_\text{in}\cdot n_\text{out} + n_\text{out}).
\]
У сети «1--8--8--1» из первого опыта $N=(8+8)+(64+8)+(8+1)=97$ параметров. Это немногим больше
сотни чисел --- такая сеть помещается в обычную переменную.

Для классификации последний слой имеет по нейрону на класс, а ответы превращают в доли функцией
\textbf{softmax}: если слой выдал $z=(2{,}0;\,1{,}0;\,0{,}1)$, то доли равны
$e^{z_k}/\sum e^{z}$, то есть примерно $(0{,}66;\,0{,}24;\,0{,}10)$. Они складываются в единицу.

\begin{caution}
Доля после softmax --- не «вероятность правоты». Сеть \emph{всегда} распределит единицу между известными
классами, даже если на вход пришло то, чего она не видела. Поэтому в классификаторе нужны
порог уверенности, класс «неизвестно» и независимая проверка.
\end{caution}

\subsection{Как сеть учится}

Обучение подбирает числа весов так, чтобы ответы на примерах были близки к правильным.
Для этого вводят \textbf{функцию потерь} --- число, показывающее, насколько ответ плох (для
синуса --- средний квадрат ошибки). Затем многократно вычисляют, в какую сторону изменить каждый вес,
чтобы потери уменьшились, и делают небольшой шаг; размер шага называют \textbf{скоростью обучения}.
Один проход по всем примерам --- \textbf{эпоха}. Вычисления ведутся на компьютере,
и вам не придётся писать их формулы: в файле \code{tinyml\_nn.py} примерно сорок строк.

Главная опасность --- \textbf{переобучение}: сеть запоминает примеры вместо общего правила.
Поэтому данные делят на части: по одной сеть учится, по другой выбирают настройки,
третью не трогают до самого конца и используют только для итоговой оценки.

\begin{figure}[H]
\centering
\begin{tikzpicture}[font=\sffamily\scriptsize,x=1cm,y=1cm]
  \draw[->] (0,0) -- (8.4,0) node[below left]{эпохи обучения};
  \draw[->] (0,0) -- (0,3.4) node[above]{ошибка};
  \draw[thick,okgreen] plot[smooth] coordinates {(0.3,3.0) (1,1.8) (2,1.1) (3.5,0.65) (5,0.4) (6.5,0.25) (8,0.15)};
  \draw[thick,warn] plot[smooth] coordinates {(0.3,3.1) (1,1.9) (2,1.3) (3.5,1.0) (5,1.1) (6.5,1.4) (8,1.8)};
  \node[okgreen,anchor=south] at (6.6,0.55) {на обучающих};
  \node[warn,anchor=south] at (7.2,1.95) {на новых};
  \draw[dashed] (3.5,0) -- (3.5,2.6) node[above]{остановиться здесь};
\end{tikzpicture}
\caption{Переобучение: ошибка на обучающих данных падает всё время, а на новых после какого-то момента растёт}
\end{figure}

\subsection{Квантование: как уместить модель}

Обычно вес --- число с плавающей запятой, 4 байта. Но для маленьких сетей точность в семь
знаков не нужна. \textbf{Квантование} хранит вес как целое число $q$ в диапазоне от $-127$ до $127$
(1 байт) и общий для слоя множитель \textbf{scale}:
\[
w \approx q\cdot\text{scale},\qquad \text{scale}=\frac{\max|w|}{127},\qquad q=\operatorname{round}\!\left(\frac{w}{\text{scale}}\right).
\]
Пример: веса $(0{,}30;\,-1{,}27;\,0{,}64)$, максимум по модулю $1{,}27$, $\text{scale}=0{,}01$,
получаем $q=(30;\,-127;\,64)$. Ошибка каждого веса не превышает $\text{scale}/2=0{,}005$.
Память весов уменьшилась в четыре раза, а сеть, как правило, почти не изменила своих ответов.

В нашем коде квантуются только веса (смещения и расчёт остаются во float), а в промышленных
инструментах, таких как TensorFlow Lite Micro, квантуются также промежуточные значения.
Тогда расчёт идёт целыми числами --- быстрее и экономнее по энергии. Для этого требуется
\emph{калибровка}: модели показывают типичные данные, чтобы подобрать диапазоны.
Но принцип остаётся тем же: меньше битов $\to$ меньше памяти и быстрее счёт, цена --- небольшая
потеря точности, \emph{которую нужно измерить} после преобразования.

\subsection{Опыт 1: сеть повторяет синус}

Начнём с задачи, у которой есть готовый ответ, ---
$y=\sin x$ на отрезке от 0 до $2\pi$. Её удобно проверять: правильное значение известно
в любой точке. Значит, качество модели можно \emph{измерить} числом.

Для опыта нужен только светодиод на \pin{4} через резистор \SI{220}{\ohm} (как в уроке 4):

\begin{figure}[H]
\centering
\begin{circuitikz}[scale=0.95,transform shape]
  \node[blk,fill=blkcpu,minimum width=2.8cm,minimum height=2.6cm] at (0,0) {};
  \node[font=\sffamily\small\bfseries] at (0,1.7) {ESP32-S3};
  \foreach \y/\n in {0.6/{\pin{4}},-0.6/GND}
    \node[font=\sffamily\scriptsize,anchor=east] at (1.4,\y) {\n};
  \draw[thick,wyellow] (1.4,0.6) -- (2.4,0.6);
  \draw[thick] (2.4,0.6) to[R,l=\SI{220}{\ohm}] (4.6,0.6) to[led,l=LED] (6.6,0.6) -- (7.4,0.6);
  \draw[thick,wblack] (1.4,-0.6) -- (7.4,-0.6) -- (7.4,0.6);
\end{circuitikz}
\caption{Подключение светодиода для опыта 1: \pin{4} через \SI{220}{\ohm} на GND}
\end{figure}

\paragraph{Этап 1. Обучение на компьютере.} В папке \code{examples/tinyml} лежат три файла:
\code{tinyml\_nn.py} (обучение, квантование, экспорт), \code{train\_sine.py} и \code{train\_gesture.py}.
Нужен Python и библиотека NumPy (\code{pip install numpy}).

\begin{Verbatim}[fontsize=\small,bgcolor=codebg]
cd examples/tinyml
python train_sine.py > sine_model.txt
\end{Verbatim}

Сообщения о качестве выводятся в консоль, а готовый блок на C++ (массивы весов, множители scale,
смещения) --- в файл. Скрипт сам себя проверяет на 400 \emph{новых} точках, не использованных при обучении:

\begin{Verbatim}[fontsize=\small,bgcolor=codebg]
макс. ошибка float32: 0.026
макс. ошибка int8:    0.044
средняя ошибка int8:  0.010
\end{Verbatim}

Видно цену квантования: максимальная ошибка выросла с 0,026 до 0,044, но на 80 байтах весов.
Для яркости светодиода, где глаз различает единицы процентов, этого достаточно; для измерительного
прибора --- нет. Решение о том, «достаточно ли», зависит от задачи, а не от модели.

\paragraph{Этап 2. Сеть в прошивке.} Содержимое сгенерированного блока вставляют между строками
\code{НАЧАЛО/КОНЕЦ СГЕНЕРИРОВАННОЙ МОДЕЛИ}. Весь расчёт сети --- это одна функция
\code{dense()} (слой), вызванная три раза. Она записана так, как слой выглядит в формуле:
для каждого нейрона суммируем «вход $\times$ вес», умножаем на scale и добавляем смещение.

Собирать этот скетч мы будем по шагам: сначала поймём, как считает один слой, затем как хранить его
веса в одном байте, и только потом подключим настоящую сеть и светодиод.

\subsubsection{Шаг 1. Слой нейронов на заготовленных числах}

Возьмём слой с двумя входами и тремя нейронами. Веса мы не обучаем, а выбираем сами --- те же числа,
что в примере выше: для входа $(0{,}5;\,-1)$ первый нейрон должен выдать $0{,}3$, а второй сработать
в ноль. Программа делает ровно то, что написано в формуле: умножить, сложить, добавить смещение, ReLU.

\codecap{Шаг 1. Один слой, веса float}
\codeinput{l42_sine_nn_s1}

\begin{itemize}[leftmargin=*]
  \item \code{W[i * N\_OUT + j]} --- вес от входа \code{i} к нейрону \code{j}. Скрипт обучения сохраняет таблицу
        «по строкам», поэтому порядок в Python и C++ совпадает; перепутаете оси --- ответы станут случайными;
  \item \code{sum = B[j]} начинает сумму со смещения, а внутренний цикл добавляет «вход $\times$ вес»;
  \item в мониторе должны появиться строки \code{0.30}, \code{0.00} и \code{1.60}: проверьте их вручную.
\end{itemize}

\subsubsection{Шаг 2. Квантование: веса в один байт}

Теперь те же шесть весов хранятся не во float (4 байта), а в \code{int8\_t} (1 байт) вместе с общим
множителем \code{scale}. Функция \code{quantize()} делает ровно то, что написано в разделе про квантование:
находит максимум по модулю, делит на 127 и округляет.

\codecap{Шаг 2. Квантование слоя в int8}
\codeinput{l42_sine_nn_s2}

\begin{itemize}[leftmargin=*]
  \item \code{lroundf(w[i] / scale)} превращает вес в целое число от $-127$ до $127$;
  \item в \code{dense()} теперь два множителя: вход на целый вес, а \code{scale} применяется один раз
        на нейрон, после суммирования;
  \item результаты немного отличаются от шага 1 (около \code{0.29} вместо \code{0.30}) --- это и есть
        цена квантования. Сравнение ответов до и после --- обязательная часть работы.
\end{itemize}

\subsubsection{Шаг 3. Настоящая сеть 1--8--8--1}

Теперь вместо шести придуманных весов --- сгенерированный блок скрипта \code{train\_sine.py}. Весь расчёт
сети --- три вызова той же \code{dense()}. В книге блок весов \emph{сокращён} многоточием, чтобы не
занимать страницу: полный блок вставляется в скетч между строками \code{НАЧАЛО/КОНЕЦ СГЕНЕРИРОВАННОЙ
МОДЕЛИ}, а готовый скетч с весами лежит в файле \code{l42\_sine\_nn.cpp}.

\codecap{Шаг 3 (фрагмент). Сеть 1--8--8--1, блок весов сокращён}
\codeinput{l42_sine_nn_s3}

\begin{itemize}[leftmargin=*]
  \item у каждого слоя свой набор: \code{W} (веса int8), \code{S} (scale) и \code{BIAS} (смещения);
  \item \code{x / (2 * PI)} --- та же \emph{нормализация} входа, что при обучении. Это самая частая причина
        ошибки: данные в прошивке подготовлены не так, как во время обучения;
  \item последний вызов \code{dense()} идёт с \code{relu = false}: выход регрессии может быть отрицательным.
\end{itemize}

\subsubsection{Шаг 4. Самопроверка: ошибка и время}

Сеть бесполезна, пока мы не знаем, насколько она точна и быстра. Функция \code{selfTest()} сравнивает
сеть со стандартной \code{sinf()} на 100 точках и усредняет время по 1000 запускам.

\codecap{Шаг 4 (фрагмент). Ошибка и время предсказания}
\codeinput{l42_sine_nn_s4}

\begin{itemize}[leftmargin=*]
  \item \code{worst} --- максимальное отличие от настоящего синуса, то есть качество, выраженное числом;
  \item \code{micros()} до и после тысячи вызовов, делённое на 1000, даёт время одного предсказания;
  \item \code{volatile float sink} не даёт компилятору выбросить цикл, результат которого нигде не нужен.
\end{itemize}

\subsubsection{Шаг 5. Светодиод повторяет синус}

Остаётся подать ответ сети на светодиод. В \code{loop()} угол \code{x} делает оборот за 4 секунды, сеть
возвращает значение от $-1$ до $1$, а мы переводим его в ШИМ $0\ldots255$ формулой $(y+1)\cdot127{,}5$ (урок 4).
Собранный из шагов 1--5 скетч с полным блоком весов --- файл \code{l42\_sine\_nn.cpp} в папке \code{code}.

\codecap{Шаг 5 (фрагмент). Яркость следует за сетью}
\codeinput{l42_sine_nn_s5}

Откройте Serial Plotter: график \code{y} должен быть похож на синусоиду с периодом 4 секунды,
а светодиод плавно «дышит». В Serial Monitor выводятся максимальная ошибка и время инференса:
запишите оба числа в тетрадь, они понадобятся в задании 2.

\subsection{Опыт 2: распознавание жестов по акселерометру}

Теперь задача, для которой нет формулы: различить три состояния платы с MPU6050 (урок 26) ---
\emph{покой}, \emph{тряска} и \emph{наклон}. Можно придумать свои классы, например «круг» и «удар»,
но их должно быть заметно меньше десяти.

Модуль подключается по I\textsuperscript{2}C так же, как в уроке 26: SDA к \pin{8}, SCL к \pin{9},
вывод AD0 на землю (адрес \code{0x68}, как в коде).

\begin{figure}[H]
\centering
\begin{circuitikz}[scale=0.95,transform shape]
  \node[blk,fill=blkcpu,minimum width=2.8cm,minimum height=4.4cm] (m) at (0,0) {};
  \node[blk,fill=blkper,minimum width=2.8cm,minimum height=4.4cm] (d) at (8,0) {};
  \node[font=\sffamily\small\bfseries] at (0,2.65) {ESP32-S3};
  \node[font=\sffamily\small\bfseries,align=center] at (8,2.85) {MPU6050\\[-2pt]\scriptsize модуль GY-521};
  \foreach \y/\n in {1.6/3V3,0.8/GND,-0.8/{\pin{9}},-1.6/{\pin{8}}}
    \node[font=\sffamily\scriptsize,anchor=east] at (1.4,\y) {\n};
  \foreach \y/\n in {1.6/VCC,0.8/GND,0/AD0,-0.8/SCL,-1.6/SDA}
    \node[font=\sffamily\scriptsize,anchor=west] at (6.6,\y) {\n};
  \draw[thick,wred]    (1.4,1.6)  -- (6.6,1.6);
  \draw[thick,wblack]  (1.4,0.8)  -- (6.6,0.8);
  \draw[thick,wblack]  (6.0,0.8)  -- (6.0,0) -- (6.6,0);
  \fill[wblack] (6.0,0.8) circle (1.6pt);
  \draw[thick,wyellow] (1.4,-0.8) -- (6.6,-0.8);
  \draw[thick,wgreen]  (1.4,-1.6) -- (6.6,-1.6);
  \node[font=\sffamily\scriptsize,text=wyellow!70!black] at (4,-0.5) {SCL};
  \node[font=\sffamily\scriptsize,text=wgreen!70!black] at (4,-1.3) {SDA};
\end{circuitikz}
\caption{Подключение MPU6050 (GY-521): питание \SI{3.3}{V}, SDA на \pin{8}, SCL на \pin{9}.
Подтяжки шины уже стоят на модуле; AD0 на GND задаёт адрес \code{0x68}}
\end{figure}

\paragraph{Этап 1. Данные.} Прошивка сбора берёт окно из 50 отсчётов (1 секунда) и сразу считает
\emph{признаки}: среднее и разброс (стандартное отклонение) по каждой из осей. Вместо 150 чисел
сети нужны всего шесть --- это уменьшает модель и ускоряет расчёт.

Скетч сбора данных тоже соберём по шагам.

\subsubsection{Шаг 1. Чтение датчика}

Сначала убедимся, что плата видит акселерометр. Датчик опрашивается 50 раз в секунду по таймеру на
\code{millis()}, а не через \code{delay()}: период остаётся ровным, даже если вывод в монитор занимает время.

\codecap{Шаг 1. Три оси ускорения в Serial Monitor}
\codeinput{l42_imu_collect_s1}

\begin{itemize}[leftmargin=*]
  \item \code{readAccel()} читает шесть байт начиная с регистра \code{0x3B} (урок 26) и переводит их в
        единицы $g$: в покое одна из осей покажет около $\pm1$;
  \item \code{nextMs += PERIOD\_MS} задаёт следующий момент опроса без накопления ошибки;
  \item если датчик не отвечает, печатается \code{sensor error}: проверьте SDA, SCL и питание.
\end{itemize}

\subsubsection{Шаг 2. Окно в одну секунду и среднее}

Сети нужны не 50 сырых отсчётов в секунду, а несколько чисел, описывающих окно. Складываем 50 отсчётов в
массив \code{samples} и, когда он заполнен, считаем первый признак --- среднее по каждой оси.

\codecap{Шаг 2. Окно из 50 отсчётов и средние}
\codeinput{l42_imu_collect_s2}

\begin{itemize}[leftmargin=*]
  \item \code{samples[count]} принимает очередной отсчёт; пока \code{++count < WINDOW}, функция возвращается;
  \item среднее показывает «как лежит плата»: наклон меняет средние по осям;
  \item после обработки окна \code{count} сбрасывается в 0 и копится новое окно.
\end{itemize}

\subsubsection{Шаг 3. Разброс и строка CSV}

Среднее не отличает тряску от покоя: при тряске оно почти такое же. Различает их \emph{разброс}
(стандартное отклонение). Теперь признаков шесть, и вывод оформлен как строка таблицы CSV.

\codecap{Шаг 3. Шесть признаков в формате CSV}
\codeinput{l42_imu_collect_s3}

\begin{itemize}[leftmargin=*]
  \item \code{sd[a] = sqrtf(sd[a] / WINDOW)} --- корень из среднего квадрата отклонений от среднего;
  \item в строке теперь есть номер сеанса и класс; пока они фиксированы нулями;
  \item заголовок с символом \code{\#} печатается один раз; скрипт обучения пропускает такие строки.
\end{itemize}

\subsubsection{Шаг 4. Класс и сеанс из монитора порта}

Остаётся научить скетч принимать команды: цифра выбирает класс, а \code{n} начинает новый сеанс.
Это итоговый скетч \code{l42\_imu\_collect}.

\codecap{Шаг 4. Сбор данных с командами}
\codeinput[\footnotesize]{l42_imu_collect}

\begin{itemize}[leftmargin=*]
  \item в начале \code{loop()} читаются все пришедшие символы: цифра меняет \code{label}, а \code{n}
        увеличивает \code{session};
  \item сообщения о смене класса и сеанса начинаются с \code{\#}, чтобы не попасть в обучение.
\end{itemize}

Откройте монитор порта. Цифрой выберите класс, затем делайте жест и наблюдайте строки вида
\code{0,2,0.0123,0.7012,0.7021,0.0210,0.0198,0.0205}. Команда \code{n} начинает новый
\textbf{сеанс}: новый запуск, новая поза платы, другой человек. Скопируйте строки в файл
\code{gestures.csv} (строки с \code{\#} скрипт пропускает).

\begin{center}\small
\begin{tabularx}{\textwidth}{@{}l X@{}}
\toprule
Правило сбора & Почему \\
\midrule
Не менее 3 сеансов на каждый класс, в каждом 40 и более окон & модель должна увидеть разброс условий, а не один раз \\
Меняйте позу платы, руку, темп & иначе сеть выучит «как лежала плата», а не «что такое тряска» \\
Включите класс «ничего» (покой, переход между жестами) & иначе любое движение станет одним из жестов \\
Записывайте отдельно сеансы для итоговой проверки & см. урок 37: соседние окна почти одинаковы \\
\bottomrule
\end{tabularx}
\end{center}

\paragraph{Этап 2. Обучение.} Скрипт делит данные \emph{по сеансам}: по последнему сеансу каждого
класса откладывается для теста и не участвует в обучении. Это тот самый честный тест из урока 37.

\begin{Verbatim}[fontsize=\small,bgcolor=codebg]
python train_gesture.py --csv gestures.csv > gesture_model.txt
\end{Verbatim}

Для проверки цепочки без платы есть режим \code{--demo} на искусственных данных. Он печатает
\textbf{матрицу ошибок}: строка --- настоящий класс, столбец --- ответ сети.

\begin{Verbatim}[fontsize=\small,bgcolor=codebg]
Тест (int8), точность 100.0%, окон: 180
[[60  0  0]
 [ 0 60  0]
 [ 0  0 60]]
\end{Verbatim}

Идеальные 100\% здесь получены на искусственных данных, где классы разведены нарочно далеко. Это
проверка программы, а не жестов. На ваших настоящих записях цифра будет ниже, и именно по
внедиагональным числам видно, \emph{какие} жесты путаются. Не расстраивайтесь, а добавляйте данные.

\paragraph{Этап 3. Инференс на плате.} В прошивке есть те же шесть признаков, нормализация
(\code{FMEAN}, \code{FSTD} --- среднее и разброс признаков из обучения), сеть 6--8--3 и softmax.

Скетч распознавания тоже собирается по шагам. Чтение датчика \code{readAccel()} и настройка в \code{setup()}
остаются теми же, что в сборе данных; весь готовый скетч --- файл \code{l42\_imu\_infer.cpp}.

\subsubsection{Шаг 1. Сеть 6--8--3 и softmax}

Сначала нужна сама функция «признаки $\to$ класс». Она нормализует шесть признаков, прогоняет их через два
слоя \code{dense()} и превращает ответы в доли функцией softmax. Блок весов в книге сокращён многоточием;
в скетче он целиком, и его создаёт \code{train\_gesture.py}.

\codecap{Шаг 1 (фрагмент). Сеть 6--8--3: классификатор}
\codeinput[\footnotesize]{l42_imu_infer_s1}

\begin{itemize}[leftmargin=*]
  \item \code{FMEAN} и \code{FSTD} --- среднее и разброс каждого признака из обучения; вход приводится к
        тому же масштабу;
  \item у последнего слоя три нейрона --- по одному на класс; \code{relu = false}, потому что дальше softmax;
  \item из каждого ответа вычитается максимум перед \code{expf()}: так число не переполняется;
  \item \code{confidence} --- доля лучшего класса после softmax. Это не вероятность правоты, а лишь
        «насколько уверенно сеть выбрала».
\end{itemize}

\subsubsection{Шаг 2. Признаки окна на плате}

Сети нужны те же шесть признаков, что при сборе данных. Мы выносим их расчёт в функцию
\code{windowFeatures()}. Если признаки на плате посчитаны иначе, чем во время обучения, сеть будет ошибаться.

\codecap{Шаг 2 (фрагмент). Признаки окна}
\codeinput{l42_imu_infer_s2}

\begin{itemize}[leftmargin=*]
  \item в первых трёх элементах --- средние, в следующих трёх --- разброс, ровно как в строке CSV;
  \item массив \code{samples} и размер окна те же, что в сборе данных.
\end{itemize}

\subsubsection{Шаг 3. Инференс и замер времени}

Соединяем всё в \code{loop()}: заполнили окно, посчитали признаки, спросили сеть и измерили время.
Пока выводится любой ответ, даже неуверенный.

\codecap{Шаг 3 (фрагмент). Окно $\to$ признаки $\to$ класс}
\codeinput{l42_imu_infer_s3}

\begin{itemize}[leftmargin=*]
  \item \code{micros()} до и после \code{classify()} показывает, сколько стоит сам расчёт сети;
  \item окно собирается секунду, поэтому новый ответ приходит раз в секунду.
\end{itemize}

\subsubsection{Шаг 4. Порог уверенности и подтверждение}

Последний шаг отличает игрушку от рабочего устройства. Заменяем \code{loop()} из шага 3 на версию с двумя
защитами. Весь скетч с этим \code{loop()} --- итоговый файл \code{l42\_imu\_infer.cpp}.

\codecap{Шаг 4 (фрагмент). Порог и подтверждение}
\codeinput[\footnotesize]{l42_imu_infer_s4}

Обратите внимание на три приёма, которые отличают игрушку от рабочего устройства:
\begin{itemize}
  \item \textbf{Порог уверенности} \code{MIN\_CONFIDENCE}: если лучший класс получил меньше 80\%,
        ответ --- «?» (неизвестно). Порог подбирают на тестовых сеансах: ниже --- больше ложных срабатываний,
        выше --- больше пропусков.
  \item \textbf{Подтверждение} \code{CONFIRM}: действие запускается, только когда три окна подряд дали
        один и тот же класс. Единичная ошибка не включит мотор (мысль из урока 37).
  \item \textbf{Измерение времени}: в каждой строке вывода --- время расчёта в микросекундах.
        Оно складывается из считывания окна (1 с) и самого инференса; на фоне секунды расчёт сети
        ничтожен. Оценивайте задержку \emph{всей} цепочки, а не одной функции.
\end{itemize}

\begin{caution}
Модель в \code{l42\_imu\_infer} обучена на демонстрационных данных и на реальной плате будет
ошибаться. Замените блок между строками \code{НАЧАЛО/КОНЕЦ СГЕНЕРИРОВАННОЙ МОДЕЛИ} результатом
вашего \code{train\_gesture.py}, иначе вы проверяете чужие данные. Если вы меняете число классов или
признаков, измените также \code{NAMES}, размеры массивов \code{in[]} и \code{h[]} и параметры вызовов \code{dense()}.
\end{caution}

\subsection{Профессиональные инструменты}

Собственный код на NumPy нужен, чтобы понять принцип: за рамками учебных моделей разумнее
использовать готовые инструменты. Они делают то же самое, но умеют сложные слои (свёртки для
изображений и звука), полное квантование и оптимизированные ядра для процессора ESP32-S3.

\begin{center}\small
\begin{tabularx}{\textwidth}{@{}>{\raggedright\arraybackslash}p{52mm} X@{}}
\toprule
Инструмент & Для чего \\
\midrule
TensorFlow Lite for Microcontrollers (TFLM) и его порт для Espressif, \code{esp-tflite-micro} & запуск обученных в Python моделей на микроконтроллере; ядра ESP-NN используют возможности ESP32-S3 \\
ESP-DL (с преобразователем ESP-PPQ) & библиотека Espressif для инференса и квантованных моделей на ESP32; нужна версия ESP-IDF, указанная в документации \\
Edge Impulse & сайт, где собирают данные, обучают и экспортируют готовую библиотеку для Arduino или ESP-IDF \\
ESP-SR, ESP-WHO & готовые решения Espressif: распознавание ключевых слов и речи, лица и объекты \\
\bottomrule
\end{tabularx}
\end{center}

Типовой маршрут с TFLM состоит из тех же шагов цикла TinyML. Модель описывают в Keras
и обучают на компьютере, затем преобразуют в формат \code{.tflite} с полным квантованием в int8:

\begin{Verbatim}[fontsize=\footnotesize,bgcolor=codebg]
converter = tf.lite.TFLiteConverter.from_keras_model(model)
converter.optimizations = [tf.lite.Optimize.DEFAULT]
converter.representative_dataset = representative_data   # типичные входы для калибровки
converter.target_spec.supported_ops = [tf.lite.OpsSet.TFLITE_BUILTINS_INT8]
converter.inference_input_type = tf.int8
converter.inference_output_type = tf.int8
open("model.tflite", "wb").write(converter.convert())
\end{Verbatim}

Файл превращают в массив C++ (командой \code{xxd -i} или коротким скриптом на Python) и подключают
к прошивке. Схема кода на плате выглядит так (имена зависят от версии библиотеки,
сверяйтесь с её примерами):

\begin{Verbatim}[fontsize=\footnotesize,bgcolor=codebg]
const tflite::Model* model = tflite::GetModel(model_data);   // массив из .tflite
static tflite::MicroMutableOpResolver<3> ops;                // только нужные операции
ops.AddFullyConnected();  ops.AddRelu();  ops.AddSoftmax();

alignas(16) static uint8_t arena[20 * 1024];                 // рабочая память модели
static tflite::MicroInterpreter interp(model, ops, arena, sizeof(arena));
interp.AllocateTensors();                                    // не хватило арены -- ошибка

interp.input(0)->data.int8[0] = quantizedInput;              // положить вход
interp.Invoke();                                             // посчитать
int8_t result = interp.output(0)->data.int8[0];              // забрать ответ
\end{Verbatim}

Здесь «арена» (\emph{tensor arena}) --- заранее выделенный кусок RAM для всех промежуточных значений.
Её размер подбирают опытом: начинают с запаса и уменьшают, пока \code{AllocateTensors()} успешен.
В PlatformIO (урок 41) такие библиотеки подключают строкой в \code{lib\_deps}. Разберитесь сначала с
приведённым выше «ручным» кодом: он делает то же, что \code{Invoke()}, только слои у вас на виду.

\begin{fact}
ESP32-S3 имеет векторные инструкции, ускоряющие операции сразу над несколькими значениями. Но ускорение получают
только оптимизированные ядра (ESP-NN, ESP-DL) и только на поддерживаемых операциях. Наш цикл на C++ из
опыта 1 выполняется обычными командами процессора. Это не недостаток, а повод измерять, а не верить рекламе.
\end{fact}

\subsection{Что измерять и как не обмануться}

\begin{center}\small
\begin{tabularx}{\textwidth}{@{}l X X@{}}
\toprule
Что & Как измерить & Зачем \\
\midrule
Размер модели & размер массивов весов в байтах (у нас 80 и 72) & хватит ли Flash и сколько останется приложению \\
Память во время работы & \code{ESP.getFreeHeap()} до и после создания модели, размер арены & модель не должна вытеснить Wi-Fi и стек \\
Время инференса & \code{micros()} до и после, усреднить по 1000 запускам (как в опыте 1) & уложится ли расчёт в период опроса датчика \\
Качество & матрица ошибок на отдельных сеансах, доля ложных срабатываний & модель лучше простого порога или нет \\
Энергия & $E=U\cdot I\cdot t$: напряжение, ток и время расчёта & хватит ли аккумулятора (урок 27) \\
\bottomrule
\end{tabularx}
\end{center}

Пример расчёта энергии с условными числами: инференс длится 2 мс при токе 80 мА и напряжении 3,3 В.
Тогда $E = 3{,}3\cdot0{,}08\cdot0{,}002 \approx 0{,}00053$ Дж, то есть 0,53 мДж на одно решение.
Если решение нужно раз в секунду, расчёт потребляет в среднем 0,53 мВт --- но датчик, радио и
стабилизатор могут расходовать на два порядка больше (урок 36). Поэтому энергию меряют у \emph{всего устройства}.

Три источника самообмана в TinyML:
\begin{itemize}
  \item \textbf{Тест на обучающих данных.} Почти любая модель «хороша» на тех примерах, на которых училась.
  \item \textbf{Данные одной обстановки.} Сеть выучила стол и руку одного человека, а не жест.
  \item \textbf{Среднее скрывает редкое.} 95\% точности могут означать, что один класс не распознаётся
        никогда. Смотрите на матрицу целиком.
\end{itemize}

Модель --- часть прошивки и должна иметь версию. Если вы меняете признаки (длину окна, ось),
старые веса теряют смысл: храните номер формата рядом с моделью, как в NVS (урок 37). Не записывайте
чужие звук и видео без разрешения: даже локальное устройство, которое «слышит», вызывает вопросы
приватности. Пока решение принимает не человек, действие должно быть безопасным при ошибке.

\begin{trouble}
\begin{itemize}[leftmargin=*]
  \item \textbf{Сеть выдаёт одно и то же число.} Проверьте нормализацию входа и порядок осей весов:
        самый частый виновник --- данные при обучении и на плате подготовлены по-разному.
  \item \textbf{Модель идеальна на компьютере и плоха на плате.} Признаки в прошивке посчитаны иначе
        (другие единицы, другая длина окна, \code{int} вместо \code{float}). Выведите признаки на плате и
        сравните с теми, что видел скрипт.
  \item \textbf{Все жесты распознаются как один класс.} Классы слишком похожи или данных мало; посмотрите
        на матрицу ошибок и графики признаков.
  \item \textbf{Ошибка компиляции в сгенерированном блоке.} Скрипт запущен с другим числом классов, чем
        ожидает прошивка (\code{NC}, \code{NAMES}, размер \code{z[]}).
  \item \textbf{Данные собираются «рывками».} Опрос датчика выполняется реже 50 Гц: избегайте длинных
        \code{delay()} и тяжёлого вывода в цикле.
\end{itemize}
\end{trouble}

\begin{summary}
\begin{itemize}[leftmargin=*]
  \item TinyML обучает модель на компьютере, а на плате остаётся только расчёт; для маленьких задач это
        десятки или сотни байт весов.
  \item Квантование уменьшает вес с 4 байт до 1 за небольшую потерю точности, которую нужно измерить.
  \item Качество проверяют на отдельных сеансах, смотрят на матрицу ошибок, а в устройстве добавляют порог
        уверенности и подтверждение, потому что softmax всегда выбирает какой-то класс.
  \item Готовые инструменты (TFLM, ESP-DL, Edge Impulse) устроены по тому же циклу, но умеют больше
        слоёв и используют ускорение процессора ESP32-S3.
\end{itemize}
\end{summary}

\begin{selfcheck}
\begin{enumerate}
  \item Что выигрывают и чем рискуют при квантовании весов из float32 в int8?
  \item Почему высокая точность на обучающих данных ничего не говорит о качестве модели?
  \item Зачем в классификаторе жестов нужны порог уверенности и несколько подтверждений подряд?
\end{enumerate}
\end{selfcheck}

\needspace{20\baselineskip}
\begin{tasks}
\begin{enumerate}
\item \lv{1} Запустите \code{train\_sine.py} и сравните ошибки float32 и int8 в консоли.
\item \lv{1} Прошейте опыт 1 и запишите максимальную ошибку и время одного предсказания.
\item \lv{2} Замените размеры слоёв на 1--4--4--1 и 1--16--16--1; сравните число весов, ошибку и время (размеры в вызовах \code{dense()} в прошивке тоже изменятся).
\item \lv{2} Научите сеть другой функции (например, $\cos x$ или $|x-\pi|$) и измерьте её качество на новых точках.
\item \lv{3} Соберите данные трёх жестов минимум в трёх сеансах на класс, обучите сеть, составьте матрицу ошибок
        на отдельном сеансе и подберите порог \code{MIN\_CONFIDENCE}.
\item \lv{4} Обучите ту же задачу в Keras, преобразуйте в int8 \code{.tflite} и запустите в TFLM;
        сравните размер, время и точность со своей реализацией.
\end{enumerate}
\end{tasks}

\paragraph{Документация для продолжения}\href{https://ai.google.dev/edge/litert/microcontrollers/overview}{TensorFlow Lite для микроконтроллеров},
\href{https://github.com/espressif/esp-tflite-micro}{esp-tflite-micro},
\href{https://github.com/espressif/esp-dl}{ESP-DL}.
